\scanpage{118}
\mspage{124}
savoir vérifier que des flèches $\unsure{\phi|q}:Z'\unsure{\mathord{-}|\to}Z$ dans Diag sont des $\unsure{\mathbb D|D}$-équiv. cohomologiques, donc \deleted{\unsure{en fait que ce}} \inserted{que}

\[H^{\unsure{\cdot|'}}_{\mathbb D}\unsure{(|[}Z,F)\xrightarrow{\sim}H^{\cdot}_{\mathbb D}(Z',\phi^*(F))\]

est iso pour tout $\unsure{F|\mathbb F}\in\operatorname{Ob}\mathbb D(Z)$ qui est loc \unsure{constante} sur $Z$. Or je ne sais vérifier un énoncé de cette nature (en dehors du cas où je sais d'avance que $Z$ est complètement $\mathbb D$-asphérique, i.e. que les coeff. loc. constants sur $Z$ sont en fait constants) \uline{que} si $\unsure{q|\phi}$ est déjà une Hot-équivalence. C'est dire que je ne sais \deleted{décrire} \inserted{exhiber} des $\mathbb D$-fibrations cohomologiques, que quand ce sont déjà des HOT-fibrations cohomologiques. Je veux, dans la suite, donner des critères utilisables pour des HOT-fibrations. Mais on va axiomatiser \deleted{\unsure{ceci}} ces développements, en gardant un $\mathbb D$ général, sur lequel je vais faire une hypothèse \inserted{si} draconienne que pratiquement, je ne connais que le dérivateur HOT qui y satisfasse. Il s'agit de \unsure{ce|ces} \unsure{ci-dessous}

\uuline{Der 3} c$_1$) Toute $\mathbb D$-équivalence est une $\mathbb D$-équiv. cohomologique

c$_2$) Toute $\mathbb D$-équivalence est une $\mathbb D$-équiv. \unsure{homologique}
% archive sheet number: 118
